\documentclass[12pt, aspectration=54]{beamer}
\usetheme{Boadilla}
\usefonttheme{serif} 
\setbeamertemplate{background} 
{
    \includegraphics[width=\paperwidth,height=.98\paperheight]{./Themes/christ_4_3_content}
}
\usepackage[utf8]{inputenc}
\usepackage[english]{babel}
% Packages needed for algorithms and pseudocode
\usepackage{algorithm, algorithmic}
%\usepackage{algpseudocode}
%\usepackage[ruled,vlined]{algorithm2e}
% Package needed for mathematical symbols and formatting
\usepackage{amsmath}
\usepackage{amsfonts}
\usepackage{amssymb}
\usepackage{ragged2e}
\usepackage{graphicx}
\graphicspath{{./Images}}
\author{Dr. S. Chanti}
%\title{Chirst Univeristy Template}
%\subtitle{Subtitle here}
\setbeamercovered{transparent} 
\institute[Assistant Professor]{Assistant Professor, \\ Department of Computer Science,\\ School of Sciences, \\ Christ (Deemed to be University),\\ Central Campus, Bengaluru} 
\date{} 
%\subject{} 
\addtobeamertemplate{frametitle}{\vspace*{0.3cm}}{\vspace*{0.2cm}}

\begin{document}
\setbeamertemplate{navigation symbols}{}

\title{Christ University Beamer Template}
\subtitle{Template 1}
{
\setbeamertemplate{background}{\includegraphics[width=\paperwidth,height=.98\paperheight]{./Themes/christ_4_3_title.pdf}}
\begin{frame}
\vspace{.5cm}
\titlepage
\end{frame}
}

\begin{frame}{Table of Contents}
    \tableofcontents
\end{frame}

\section{Introduction}
\begin{frame}{Introduction}
\justifying
    Beamer is a LaTeX  document class for creating presentation slides, with a wide range of templates and a set of features for making slideshow effects.
\end{frame}

\section{Different Types of Blocks in Latex Beamer}
\begin{frame}{Different Types of Blocks in Latex Beamer}
    \begin{block}{Using Block Command}
       The block command will help you in highlighting the content. 
    \end{block}
    \begin{alertblock}{Using Alert Block Command}
        Alertblock can be used to highlight very important points
    \end{alertblock}
    \begin{example}
        This block can be used to provide examples.
    \end{example}
\end{frame}

\section{Inserting an Image}
\begin{frame}{Inserting an Image}
\centering
\includegraphics[width=0.5\textwidth]{Images/Christ_logo_transparent.jpg}
    
\end{frame}
\begin{frame}{Images with Different Resolution and Types}
    \begin{columns}
\begin{column}{0.5\textwidth}
   \begin{center}
     \includegraphics[width=0.85\textwidth]{Images/2a.png}      
     \end{center}
\end{column}
\begin{column}{0.5\textwidth}
    \begin{center}
     \includegraphics[width=0.85\textwidth]{Images/2b.png}      
     \end{center}
\end{column}
\end{columns}
\end{frame}

\begin{frame}{Image in PDF Format}
    \begin{center}
     \includegraphics[width=0.75\textwidth]{Images/2c.pdf}      
     \end{center}
\end{frame}

\section{Introducing Two Columns}
\begin{frame}{Introducing Two Columns}
\begin{columns}
\begin{column}{0.5\textwidth}
   \begin{block}{Mission}
       CHRIST (Deemed to be University) is a nurturing ground for an individual's holistic development to make effective contribution to the society in a dynamic environment.
   \end{block}
\end{column}
\begin{column}{0.5\textwidth}
    \begin{center}
     \includegraphics[width=0.85\textwidth]{Images/Christ_logo.jpg}      
     \end{center}
\end{column}
\end{columns}
\end{frame}

\section{Inserting a Table}
\begin{frame}{Inserting a Table}
\begin{table}[]
\begin{tabular}{|c|l|}
\hline
\textbf{S. No.} & \multicolumn{1}{c|}{\textbf{Core Values}} \\ \hline
1 & Faith in God \\ \hline
2 & Moral Uprightness \\ \hline
3 & Love of Fellow Beings \\ \hline
4 & Social Responsibility \\ \hline
5 & Pursuit of Excellence \\ \hline
\end{tabular}
\end{table}
\end{frame}

\section{Mathematical Equations}
\begin{frame}{Mathematical Equations}
    \begin{block}{Formula}
    calculated in this equation
    \begin{equation}    % <--- deleted empty lines
            Total = \frac{1}{n}*\sum_{i=0}^n (eachnode_i)
    \end{equation}
    \begin{equation}    % <--- deleted empty lines
            Total = \frac{1}{n}*\sum_{i=0}^n (eachnode_i)
    \end{equation}
\end{block}
\end{frame}

\section{Inserting Bullet Points}
\begin{frame}{Inserting Bullet Points}
    \begin{itemize}
        \item This is how we define bullet points
        \item[*] We can the bullet points according our needs.
        \item[\$] Let's try with some different symbols.
    \end{itemize}
\vspace{.5cm}
    \begin{enumerate}
        \item In case if you need numbers, you can try \textit{enumerate} option.
        \item In case if you need numbers, you can try \textit{enumerate} option.
    \end{enumerate}
\end{frame}

\section{Algorithm(s) in Latex Beamer}
\begin{frame}{Writing an Algorithm}
\begin{algorithm}[H]
\begin{algorithmic}[1]
\FOR{$i=1$ to $N$}
\FOR{$j=1$ to $JJJJ$}
\STATE $energy[i*JJJ+j] =$ 
$ interpolate(AAA[i*JJJ+j], ZZZ)$
\ENDFOR
\ENDFOR
\end{algorithmic}
\caption{pseudocode for the calculation of }
\label{alg:seq}
\end{algorithm}
\end{frame}
\end{document}
